\section*{Введение}
\addcontentsline{toc}{section}{Введение}

В современных условиях эффективное управление личными и семейными финансами является важной задачей для большинства людей. Однако многие семьи сталкиваются с трудностями при контроле доходов и расходов.

\textbf{Проблема.} Основная проблема состоит в том, что у людей часто нет удобного и надёжного инструмента для систематического учёта финансов. Многие ведут записи в заметках или простых таблицах Excel. Это приводит к неконтролируемым тратам, невозможности правильно спланировать бюджет и снижению финансовой стабильности. Кроме того, существующие приложения часто неудобны для совместного использования несколькими людьми, имеют ограниченные возможности аналитики и не всегда позволяют работать с чеками и документами.

\textbf{Актуальность.} Тема учёта семейных финансов остаётся актуальной в 2026 году. Согласно исследованию Аналитического центра НАФИ и страховой компании «Росгосстрах Жизнь» (2025)~\cite{nafi2025}, индекс финансовой грамотности россиян снизился до 12,61 балла из 21 возможного. Доля граждан с низким уровнем финансовой грамотности выросла до 34~\%, а среди молодёжи 18--34 лет составляет 51~\%. Почти половина работающих россиян регулярно испытывает финансовые трудности.

В условиях прогнозируемой инфляции 4,5--5,5~\% (по данным Банка России на 2026 год)~\cite{cbr2026} семьи всё больше нуждаются в удобном инструменте для ведения общего семейного бюджета, отслеживания расходов всех членов семьи, загрузки чеков и получения наглядной статистики.

\textbf{Степень разработанности проблемы.} На сегодняшний день существует множество приложений для учёта финансов. Они позволяют вести учёт транзакций, ставить бюджеты и смотреть простые отчёты. Однако большинство из них имеют существенные недостатки: платные премиум-функции, слабая поддержка совместной работы нескольких пользователей, ограниченные возможности загрузки файлов и не всегда удобная аналитика. Готовых открытых веб-решений, которые одновременно включают систему ролей пользователей, полноценный CRUD-интерфейс, загрузку и выгрузку файлов чеков и глубокую агрегацию данных, относительно немного. Это определяет место и актуальность настоящей курсовой работы.

\textbf{Объект исследования} — процесс учёта и анализа доходов и расходов семейного бюджета с помощью веб-технологий.

\textbf{Предмет исследования} — методы и средства разработки веб-приложения для управления личными финансами, включающего аутентификацию, авторизацию по ролям, CRUD-операции, агрегацию данных и работу с файлами.

\textbf{Цель работы} — разработать веб-приложение для учёта расходов и доходов семейного бюджета, которое обеспечит удобный совместный контроль финансов, аналитику и безопасность данных.

Для достижения поставленной цели необходимо решить следующие \textbf{задачи}:

\begin{itemize}
    \item проанализировать существующие программные продукты и технологии разработки веб-приложений для управления финансами;
    \item определить требования целевой аудитории и спроектировать функциональность приложения (диаграммы вариантов использования, User Story);
    \item разработать модель данных (ER-диаграмму, логическую и физическую схемы базы данных);
    \item спроектировать пользовательский интерфейс (wireframes);
    \item реализовать систему аутентификации и авторизации с проверкой прав доступа по ролям;
    \item разработать CRUD-интерфейс для основных сущностей (транзакции, категории, бюджеты);
    \item реализовать механизмы загрузки и выгрузки файлов (чеки, отчёты) и агрегацию данных с визуализацией статистики;
    \item провести тестирование, деплой приложения и оформить отчётную документацию.
\end{itemize}

\textbf{Актуальность темы.} В работе применялись методы системного анализа, сравнительный анализ существующих решений, структурное и объектно-ориентированное проектирование. Для реализации использовались современные веб-технологии (backend, frontend и база данных). Визуализация данных выполнялась с помощью специализированных библиотек. Тестирование проводилось по функциональным и нефункциональным требованиям.

Структура отчёта соответствует логике поставленных задач и включает анализ предметной области, проектирование, описание реализации и заключение.
